\section{Conclusion}
\label{sec:conclusion}

We introduced DVA-MCTS, the first algorithm for test-time compute scaling with
provably efficient dynamic verifier allocation. By formulating the problem as online
resource allocation and exploiting the Lipschitz structure of process reward models,
we proved that DVA-MCTS achieves $O(\sqrt{T\log K})$ regret against the oracle
allocation policy while invoking the verifier only $O(\sqrt{T}\log T)$ times---a
sub-linear reduction in verification overhead. A matching $\Omega(\sqrt{T})$ lower
bound establishes rate-optimality.

\paragraph{Limitations.}
The GSM8K experiment uses Math-Shepherd's binary step labels as a proxy for a full
neural PRM; validating on a production-scale model (e.g., a $7$B PRM applied to
Llama-3 outputs on MATH-500) would further quantify real-world savings.
The regret definition (Definition~\ref{def:regret}) measures quality relative to
states \emph{visited} by the search; a stronger benchmark would compare against
the best leaf in the full tree, which is generally inaccessible.

\paragraph{Future directions.}
Several extensions remain:
(1) \emph{Adaptive tree depth:} the exploration--exploitation crossover occurs at
$T \approx K^D$; choosing $D$ as a function of $B$ shifts the efficiency gains
into practically relevant budget ranges.
(2) \emph{Cost-aware allocation:} when verifier cost is non-uniform (e.g., depends on
sequence length), incorporate a cost-weighted threshold.
(3) \emph{Production-scale PRM validation:} extending the GSM8K result
(Section~\ref{sec:gsm8k}) to neural PRMs on MATH-500 or AIME would quantify
real-world savings at deployment scale.
(4) \emph{Integration with speculative decoding:} combine adaptive verification with
draft-model generation for joint compute reduction.

\paragraph{Broader impact.}
Reducing the compute cost of test-time scaling has clear positive implications:
lower inference cost makes capable LLMs accessible to a broader range of users and
reduces energy consumption. We are unaware of negative societal impacts specific to
this efficiency contribution.
